problem:  
Let \( (M^m,g) \) be a closed Riemannian manifold of fixed dimension \(m\) with uniformly bounded geometry, i.e.  
\[
|\mathrm{Rm}_g| \le K_M, \qquad \mathrm{inj}(M,g) \ge i_0 > 0,
\]  
and let \( (N,h) \) be a compact Riemannian manifold with nonpositive sectional curvature and bounded geometry constants \( |\mathrm{Rm}_h| \le K_N \), \( \mathrm{inj}(N,h) \ge j_0 > 0 \).  
For each small \( \varepsilon > 0 \), consider an \( \varepsilon \)-approximately harmonic map \( u_\varepsilon : M \to N \) satisfying  
\[
\|\tau(u_\varepsilon)\|_{L^2(M)} \le \varepsilon
\quad \text{and} \quad E(u_\varepsilon) \le E_0,
\]
where \( \tau(u_\varepsilon) \) is the tension field and \(E(u_\varepsilon)\) is the energy.

We know from classical results that when \( \varepsilon \) is small compared to geometry-dependent thresholds, one has closeness of \(u_\varepsilon\) to a genuine harmonic map up to bubbling. However, the dependence of such quantitative control on geometric constants is not clearly understood.

**Problem:**  
For fixed dimension \(m\) and bounded geometry class \(\mathcal{G}(m,K_M,i_0,K_N,j_0)\), does there exist a *dimension-only* modulus of stability  
\[
\Phi_m:(0,\infty) \to (0,\infty) \text{ with } \Phi_m(t) \to 0 \text{ as } t \to 0
\]
such that for all \((M,g),(N,h)\in\mathcal{G}\) and all approximately harmonic maps \(u_\varepsilon\) as above,  
\[
\mathrm{dist}_{W^{1,2}}(u_\varepsilon,\mathcal{H}(M,N)) \le \Phi_m(\|\tau(u_\varepsilon)\|_{L^2(M)})
\]
holds, where \(\mathcal{H}(M,N)\) is the set of harmonic maps \(M \to N\)?  
Equivalently, can one establish a *dimension-dependent quantitative rigidity theorem* for the class of approximate harmonic maps between bounded-geometry manifolds?

Why is it a "good" problem:  
This revised formulation asks for a *geometrically stable*, quantitative version of harmonic map compactness — one that depends only on dimension, not on specific curvature or injectivity bounds once these are kept under control. It targets the heart of geometric analysis: how analytic quantities (tension field) capture geometric rigidity (closeness to harmonicity).  
The problem bridges nonlinear PDE analysis, geometric measure theory (through bubble-tree convergence and defect measures), and Riemannian geometry (through bounded-geometry dependence).  

Solving it would unify local ε-regularity and global compactness principles into a single quantitative stability modulus, providing a new analytical invariant of dimension. Even partial progress would likely require developing new scale-invariant estimates or a metric-analytic theory of quantitative bubble dynamics, yielding broader impact beyond harmonic maps themselves.  
Its simplicity of statement conceals profound depth, making it both accessible in form and challenging at the highest research level — meeting the ideal of a "good" mathematical problem.